% Exact TikZ call from pentagon_v005.tex; no geometric modification.
\tikzfig{Gamma}{Ice-ray dimer deformation.  The ice-ray
  dimer (center) admits a shear motion that preserves all edges of
  contact.  This deformation increases the area of the dimer.  It is
  obvious by the symmetry of the figures on the left and right that
  the area function along this deformation has a critical point at the
  ice-ray dimer.  It is known that the ice-ray dimer is a local
  minimum of the area function \cite{Kus}.  }
{
\begin{scope}[scale=0.5]
\threepent{0.00}{0.00}{72.00}{0.96}{1.43}{252.00}{1.81}{0.00}{216.00};
\threepent{0.00}{0.00}{-72.00}{0.96}{-1.43}{-252.00}{1.81}{0.00}{-216.00};
\draw (0.5,-3.5) node {$\area=2\acrit$};
\end{scope}
\begin{scope}[scale=0.5,xshift=6cm]
\threepent{0.00}{0.00}{79.49}{0.93}{1.37}{259.49}{1.82}{0.00}{223.49};
\threepent{0.00}{0.00}{-64.51}{1.03}{-1.50}{-244.51}{1.82}{0.00}{-208.51};
\draw (0.5,-3.5) node {$\area\approx 2\acrit+0.03$};
\end{scope}
\begin{scope}[scale=0.5,xshift=-6cm]
\threepent{0.00}{0.00}{-79.49}{0.93}{-1.37}{-259.49}{1.82}{0.00}{-223.49};
\threepent{0.00}{0.00}{64.51}{1.03}{1.50}{244.51}{1.82}{0.00}{208.51};
\draw (0.5,-3.5) node {$\area\approx 2\acrit+0.03$};
\end{scope}
}
